% Lösungen Einheit 2 von 4 – Ableitung an einer Stelle (zusammenbau v0.1)
\einheitenkopf{Einheit 2 von 4 · Ableitung an einer Stelle}

\erg{12}{a) Ableitungswert – Steigung oder Rate: nach 4 Jahren wächst der Baum momentan mit $35$ cm je Jahr \quad b) $h'(t) = -10t + 20$, $h'(1) = 10$ m je s – die momentane Änderungsrate bei $t = 1$ \quad c) Steigungsdreieck: 2 nach rechts, 6 nach oben, also Steigung $3$; die Ableitung einer linearen Funktion ist an jeder Stelle ihre Steigung, also $g'(3) = 3$ \quad d) $g'(x) = 3\mathrm{e}^x$, $g'(0) = 3$; Tangente im Punkt $(0 | 0)$ mit Steigungsdreieck 1 nach rechts, 3 nach oben \quad e) $p'(t) = 4t - 40 = -16$, also $t = 6$: sechs Stunden nach Messbeginn \quad f) Um 9:00 Uhr nimmt die Staulänge ab (die Rate ist negativ) – nicht: die Staulänge ist negativ. \quad g) $(3 | 180)$: Drei Stunden nach Beobachtungsbeginn nimmt der Tankinhalt momentan mit $180$ kg je Stunde zu – $180$ ist eine Rate, kein Tankinhalt. \quad h) $f'(x) = 0{,}3x^2 - 2x + 2$, $f'(0) = 2$; $p'(x) = -2x + 3$, $p'(0{,}5) = 2$ – gleiche Steigung: die Tangenten an beide Begrenzungslinien in diesen Punkten sind parallel, die Schaufel läuft dort ohne Knick weiter. \quad i) Falsch: Der Graph ist an der Wendestelle bei $x = 1$ am steilsten; die Tangente dort hat die Steigung $m = 3$ (1 nach rechts, 3 nach oben), also ist die Steigung nicht überall kleiner als $2{,}5$.}
\erg{13}{a) Tangente im Punkt $(2 | 40)$ mit Steigungsdreieck: etwa $20$ m³ je Stunde (die Steigung, nicht der Wert $40$)}
\erg{14}{a) Graph I: steigt immer langsamer (rechtsgekrümmt) – je höher der Wasserstand, desto breiter die Schale, bei gleichem Zufluss steigt die Höhe je Minute weniger}
\erg{15}{a) Fehler: Paul berechnet den Funktionswert statt des Ableitungswerts. Richtig: $g'(x) = 3\mathrm{e}^x$, $g'(0) = 3$ \quad b) Die Ableitung ist die Tangentensteigung; die Tangente an eine Gerade ist die Gerade selbst, und deren Steigung ist überall $m$ – auch jede Sekantensteigung ist $m$. \quad c) $h'(t) = -t + 12 = 4$, also $t = 8$: ab acht Minuten nach dem Start (davor steigt er schneller, $h'(0) = 12$) \quad d) Tangente durch $(2 | 2)$ mit Steigung $2$: Steigungsdreieck 1 nach rechts, 2 nach oben}

15d)

\begin{ksys}[xmin=-1,xmax=4,ymin=-1,ymax=6,klein]\parabel{0.5}{0}{0}{f}\tangentean{0.5*\x^2}{2}{t}\end{ksys}

